% Generated by 03_character_attention.py -- do not edit by hand.
% Requires \usepackage{booktabs}.
\begin{table}[t]
\centering
\begin{tabular}{lccc}
\toprule
Shot distance & M (\%) & W (\%) & W $-$ M (\%) \\
\midrule
Extreme close-up & 0.2 (0.1, 0.3) & 0.1 (0.1, 0.1) & $-$0.1 ($-$0.2, 0.0) \\
Close-up & 0.7 (0.5, 0.9) & 0.8 (0.7, 1.1) & 0.2 ($-$0.1, 0.4) \\
Medium close-up & 18.0 (16.2, 19.9) & 17.5 (15.6, 19.5) & $-$0.5 ($-$2.3, 1.5) \\
Medium & 48.6 (46.9, 50.2) & 49.0 (47.1, 50.9) & 0.4 ($-$1.5, 2.6) \\
Medium long & 9.0 (8.1, 9.9) & 10.3 (9.3, 11.4) & 1.2 (0.2, 2.4) \\
Long & 23.1 (20.9, 25.4) & 21.8 (19.9, 23.9) & $-$1.3 ($-$3.6, 0.6) \\
Extreme long & 0.4 (0.3, 0.6) & 0.4 (0.3, 0.7) & $-$0.0 ($-$0.1, 0.1) \\
\bottomrule
\end{tabular}
\caption{Share of a character's on-screen frames at each shot distance, for men (M) and women (W) among animated human characters in the year-matched animated/live-action sample, with 95\% bootstrap confidence intervals. Each movie contributes one observation and resampling is at the level of movies. All three columns are computed from the same resamples: the W $-$ M point estimate is exactly the difference of the M and W point estimates, but its interval is the percentile interval of the paired per-resample difference, not the difference of the M and W intervals. Because M and W move together from resample to resample, it is the narrower of the two and is the one to read for whether the gap differs from zero. Restricted to the 132 movies with recognized characters of both genders.}
\label{tab:gender-rate-paired-animated-human}
\end{table}
